\documentclass[11pt]{article}
\usepackage[T1]{fontenc}
\usepackage{amsmath, amssymb, amsthm}
\usepackage{mathtools}

\newtheorem{definition}{Definition}
\newtheorem{theorem}{Theorem}

\title{Bewegungsontologie}
\author{}
\date{}

\begin{document}
\maketitle

\section*{Grundbegriffe}

\paragraph{Sorten.}
Es gibt zwei Grundbereiche (Sorten):
\begin{itemize}
  \item $S$ -- die Sukzessionen, frei erzeugt durch die Konstante $s_1$ und die einstellige Nachfolgerfunktion $N : S \to S$;
  \item $B$ -- die Bewegungen, ein primitiver Bereich.
\end{itemize}

\paragraph{Primitives Prädikat.}
$$E \subseteq B \times S, \qquad E(\delta,s) \;:\; \text{``}\delta\text{ existiert in } s\text{''}$$

\paragraph{Ordnungsrelation $\prec$ auf $S$ (rekursiv definiert).}
\begin{align}
\prec(s,\, N(s')) &\;:\Longleftrightarrow\; (s = s') \vee \prec(s,s') \\
\prec(s,\, s_1)    &\;:\Longleftrightarrow\; \bot
\end{align}

\section*{Definitionen}

\begin{definition}[Notwendigkeit]
$$Nw(\delta,s) \;:\Longleftrightarrow\; E(\delta,s)$$
\end{definition}

\begin{definition}[Unmöglichkeit]
$$Um(\delta,s) \;:\Longleftrightarrow\; \neg\, E(\delta,s)$$
\end{definition}

\begin{definition}[Möglichkeit]
$$Mo(\delta,s) \;:\Longleftrightarrow\; E(\delta,s)$$
\end{definition}

\begin{definition}[Energie]
$$En(\delta,s) \;:\Longleftrightarrow\; \neg\, Mo(\delta,s) \;\wedge\; \exists s'.\big(\prec(s,s') \wedge Nw(\delta,s')\big)$$
\end{definition}

\section*{Theorem}

\begin{theorem}
$$\forall \delta \, \forall s.\; E(\delta,s) \rightarrow Nw(\delta,s)$$
\end{theorem}

\begin{proof}
Sei $\delta \in B$ und $s \in S$ mit $E(\delta,s)$. Nach Definition von $Nw$ gilt
$$Nw(\delta,s) \Longleftrightarrow E(\delta,s),$$
also folgt $Nw(\delta,s)$ unmittelbar aus der Voraussetzung. $\blacksquare$
\end{proof}

\end{document}
